\begin{proof}[Proof of \cref{obj:prop:phase-boundaries}]
\begin{enumerate}
\item\label{proof:phase-boundaries:comparison} For every \(t\in\mathbb N\), write
\[
\ell_t:=\log(e+N_t),\qquad
r_t:=r_{n_t,d_t,q_t}
=\min\left\{1,N_t^{-1}
+\left(\frac{d_t}{N_t\ell_t}\right)^2\right\},
\qquad
\mathfrak R_t:=\mathfrak R_{n_t,d_t,q_t}.
\]
The expression for \(r_t\) is supplied by \cref{obj:def:frontier-rate}. The size and arrival assumptions give
\[
N_t=n_tq_t>0,\qquad
N_t^{-1}+\left(\frac{d_t}{N_t\ell_t}\right)^2\ge0,
\qquad r_t\ge0.
\]
All hypotheses of \cref{obj:thm:matched-minimax-frontier} hold at every \(t\). Consequently, there are constants
\(0<\underline c\le\overline C<\infty\), independent of \(t\), such that
\[
\underline c\,r_t\le\mathfrak R_t\le\overline C\,r_t.
\]
In particular,
\[
0\le r_t\le\underline c^{-1}\mathfrak R_t,
\qquad
0\le\mathfrak R_t\le\overline C\,r_t.
\]
The first inequality and the squeeze theorem show that
\(\mathfrak R_t\to0\) implies \(r_t\to0\); the second gives the converse. Thus
\[
\mathfrak R_t\longrightarrow0
\quad\Longleftrightarrow\quad
r_t\longrightarrow0.
\]
% lean: CausalSmith.Stat.MarRareqLogfrontier.phase_boundaries

\item\label{proof:phase-boundaries:convergence} Define, for every \(t\in\mathbb N\),
\[
\nu_t:=N_t^{-1},
\qquad
\xi_t:=\frac{d_t}{N_t\ell_t}.
\]
Since \(e>1\) and \(N_t>0\),
\[
e+N_t>1,\qquad
\ell_t>0,\qquad
N_t\ell_t>0,\qquad
\nu_t\ge0,\qquad \xi_t\ge0.
\]
Effective sample growth gives
\[
\nu_t\longrightarrow0.
\]
Continuity of the square root, together with
\[
\sqrt{\xi_t^2}=|\xi_t|=\xi_t,
\]
shows that \(\xi_t^2\to0\) implies \(\xi_t\to0\). Continuity of squaring gives the reverse implication. Hence
\[
\xi_t^2\longrightarrow0
\quad\Longleftrightarrow\quad
\xi_t\longrightarrow0.
\]
Because \(N_t\ell_t>0\) for every \(t\), the quotient characterization of little \(o\) gives
\[
d_t=o(N_t\ell_t)
\quad\Longleftrightarrow\quad
\xi_t\longrightarrow0.
\]

We next establish the truncation equivalence
\[
r_t=\min\{1,\nu_t+\xi_t^2\}\longrightarrow0
\quad\Longleftrightarrow\quad
\xi_t^2\longrightarrow0.
\]
For the forward direction, \(r_t\to0\) implies \(r_t<1\) eventually. At each such index, \(\nu_t+\xi_t^2\ge1\) would give \(r_t=1\). Therefore, eventually,
\[
\nu_t+\xi_t^2<1,\qquad
r_t=\nu_t+\xi_t^2.
\]
The sum consequently tends to zero, and
\[
0\le\xi_t^2\le\nu_t+\xi_t^2
\]
implies \(\xi_t^2\to0\). Conversely, if \(\xi_t^2\to0\), then
\[
\nu_t+\xi_t^2\longrightarrow0,
\qquad
\min\{1,\nu_t+\xi_t^2\}\longrightarrow\min\{1,0\}=0,
\]
by continuity of addition and minimum. Combining these equivalences with \cref{proof:phase-boundaries:comparison} yields
\[
\mathfrak R_t\longrightarrow0
\quad\Longleftrightarrow\quad
d_t=o(N_t\ell_t)
=o\!\left(N_t\log(e+N_t)\right).
\]
% lean: CausalSmith.Stat.MarRareqLogfrontier.phase_min_rate_tendsto
% lean: CausalSmith.Stat.MarRareqLogfrontier.phase_frontier_rate_tendsto

\item\label{proof:phase-boundaries:theta-comparison} The comparison in \cref{proof:phase-boundaries:comparison} also gives a two-sided asymptotic comparison. Indeed, nonnegativity allows us to write, for every \(t\),
\[
|\mathfrak R_t|
\le\overline C\,|r_t|,
\qquad
|r_t|
\le\underline c^{-1}|\mathfrak R_t|.
\]
These are the defining bounds for
\[
\mathfrak R_t=O(r_t),
\qquad
r_t=O(\mathfrak R_t),
\qquad\text{and hence}\qquad
\mathfrak R_t=\Theta(r_t).
\]
% lean: CausalSmith.Stat.MarRareqLogfrontier.phase_boundaries

\item\label{proof:phase-boundaries:baseline} Since \(N_t\to\infty\), eventually \(N_t\ge1\). At those indices,
\[
0\le\nu_t\le1,
\qquad
\nu_t\le\nu_t+\xi_t^2.
\]
Thus both entries in the minimum defining \(r_t\) are at least \(\nu_t\), and
\[
0\le\nu_t\le\min\{1,\nu_t+\xi_t^2\}=r_t.
\]
This proves the baseline bound
\[
\nu_t=O(r_t).
\]
% lean: CausalSmith.Stat.MarRareqLogfrontier.phase_rate_lower_bigO

\item\label{proof:phase-boundaries:dimension-necessity} Define the dimension scale, for every \(t\in\mathbb N\), by
\[
B_{\mathrm{phase},t}:=\sqrt{N_t}\,\ell_t.
\]
The positivity established in \cref{proof:phase-boundaries:convergence} gives
\[
B_{\mathrm{phase},t}>0,
\qquad
N_t\ell_t>0.
\]
Moreover, using \((\sqrt{N_t})^2=N_t\),
\[
\nu_t(N_t\ell_t)^2
=N_t^{-1}N_t^2\ell_t^2
=N_t\ell_t^2
=B_{\mathrm{phase},t}^2.
\]

Suppose first that \(r_t=O(\nu_t)\). Choose a real constant
\(\beta_{\mathrm{rate}}\) such that, eventually,
\[
0\le r_t\le\beta_{\mathrm{rate}}\nu_t.
\]
Since \(\beta_{\mathrm{rate}}\nu_t\to0\), the squeeze theorem gives \(r_t\to0\), and hence \(r_t<1\) eventually. As in \cref{proof:phase-boundaries:convergence}, this forces
\[
\nu_t+\xi_t^2<1,
\qquad
r_t=\nu_t+\xi_t^2
\]
eventually. Intersecting these eventual conditions with the assumed bound gives
\[
\nu_t+\xi_t^2\le\beta_{\mathrm{rate}}\nu_t,
\qquad
\xi_t^2\le\beta_{\mathrm{rate}}\nu_t,
\]
where the latter inequality uses \(\nu_t\ge0\). Multiplication by the positive quantity \((N_t\ell_t)^2\) then yields
\[
d_t^2
=\xi_t^2(N_t\ell_t)^2
\le\beta_{\mathrm{rate}}\nu_t(N_t\ell_t)^2
=\beta_{\mathrm{rate}}B_{\mathrm{phase},t}^2.
\]
Set
\[
\beta_{\mathrm{bound}}:=\max\{1,\beta_{\mathrm{rate}}\}\in\mathbb R.
\]
Then
\[
\beta_{\mathrm{bound}}\ge1,\qquad
\beta_{\mathrm{rate}}\le\beta_{\mathrm{bound}}
\le\beta_{\mathrm{bound}}^2.
\]
Consequently, eventually,
\[
d_t^2
\le\beta_{\mathrm{bound}}^2B_{\mathrm{phase},t}^2.
\]
Both \(d_t\) and \(\beta_{\mathrm{bound}}B_{\mathrm{phase},t}\) are nonnegative, so this squared inequality gives
\[
d_t\le\beta_{\mathrm{bound}}B_{\mathrm{phase},t}.
\]
Therefore
\[
r_t=O(\nu_t)
\quad\Longrightarrow\quad
d_t=O(B_{\mathrm{phase},t}).
\]
% lean: CausalSmith.Stat.MarRareqLogfrontier.phase_dimension_rate_bigO

\item\label{proof:phase-boundaries:dimension-sufficiency} Conversely, suppose that \(d_t=O(B_{\mathrm{phase},t})\). Choose a real constant
\(\beta_{\mathrm{dim}}\) such that, eventually,
\[
0\le d_t\le\beta_{\mathrm{dim}}B_{\mathrm{phase},t}.
\]
The right-hand side is consequently nonnegative, so squaring gives
\[
d_t^2\le\beta_{\mathrm{dim}}^2B_{\mathrm{phase},t}^2.
\]
Dividing by \((N_t\ell_t)^2>0\) and using the identity from \cref{proof:phase-boundaries:dimension-necessity} yields
\[
\xi_t^2
=\frac{d_t^2}{(N_t\ell_t)^2}
\le
\beta_{\mathrm{dim}}^2
\frac{B_{\mathrm{phase},t}^2}{(N_t\ell_t)^2}
=\beta_{\mathrm{dim}}^2\nu_t.
\]
It follows that, eventually,
\[
0\le r_t
=\min\{1,\nu_t+\xi_t^2\}
\le\nu_t+\xi_t^2
\le(1+\beta_{\mathrm{dim}}^2)\nu_t.
\]
Thus \(r_t=O(\nu_t)\). Together with \cref{proof:phase-boundaries:dimension-necessity}, this establishes
\[
r_t=O(\nu_t)
\quad\Longleftrightarrow\quad
d_t=O(B_{\mathrm{phase},t})
=O\!\left(\sqrt{N_t}\log(e+N_t)\right).
\]
% lean: CausalSmith.Stat.MarRareqLogfrontier.phase_dimension_rate_bigO

\item If \(\mathfrak R_t=\Theta(\nu_t)\), then
\[
r_t=O(\mathfrak R_t),
\qquad
\mathfrak R_t=O(\nu_t)
\quad\Longrightarrow\quad
r_t=O(\nu_t),
\]
by transitivity of big \(O\). The equivalence in \cref{proof:phase-boundaries:dimension-sufficiency} therefore gives
\(d_t=O(B_{\mathrm{phase},t})\).

Conversely, if \(d_t=O(B_{\mathrm{phase},t})\), then \cref{proof:phase-boundaries:dimension-sufficiency} gives \(r_t=O(\nu_t)\). Combining this with \cref{proof:phase-boundaries:theta-comparison,proof:phase-boundaries:baseline} gives
\[
\mathfrak R_t=O(r_t),\quad r_t=O(\nu_t)
\quad\Longrightarrow\quad
\mathfrak R_t=O(\nu_t),
\]
and
\[
\nu_t=O(r_t),\quad r_t=O(\mathfrak R_t)
\quad\Longrightarrow\quad
\nu_t=O(\mathfrak R_t).
\]
These two bounds prove
\[
\mathfrak R_t=\Theta(N_t^{-1})
\quad\Longleftrightarrow\quad
d_t=O\!\left(\sqrt{N_t}\log(e+N_t)\right),
\]
which is the second claimed equivalence.
% lean: CausalSmith.Stat.MarRareqLogfrontier.phase_boundaries
\end{enumerate}
\end{proof}
